%------------------------------------------------------------------------------%
\documentclass[fleqn,utf8,aspectratio=169,12pt]{beamer}

\usepackage{gaal}

\title{Produto Vetorial}

%------------------------------------------------------------------------------%
\begin{document}

\addtitle

%------------------------------------------------------------------------------%
\section{Introdução}

%------------------------------------------------------------------------------%
\begin{frame}{Produto Vetorial}{}
  \onslide<1->{\emph{\color<2>{red}Apenas para vetores em $\R^3$}}

  \bigskip
  \onslide<3->{é possível definir uma operação de produto entre vetores}

  \bigskip
  \onslide<4->{que produz um novo vetor ortogonal a ambos os vetores originais}
\end{frame}

%------------------------------------------------------------------------------%
\begin{frame}{Produto Vetorial}
  Dados dois vetores não nulos
  \[
    \vec{v},\; \vec{w} \in \emph{\R^3}
  \]
  \pause
  O \emph{produto escalar} produz um número
  \[
    \vec{v} \cdot \vec{w} \in \emph{\R}
  \]
  \pause
  O \emph{produto vetorial} produz um vetor ortogonal aos dois vetores
  \[
    \vec{v} \times \vec{w} \in \emph{\R^3}
  \]
  \pause
  \emph{Não misture as notações: \; $\cdot$ \; $\times$}
\end{frame}

%------------------------------------------------------------------------------%
\begin{frame}{Definição Geométrica}
  Considere dois vetores com a mesma origem
  \[
    \vec{v},\; \vec{w} \in \R^3
  \]
  \pause
  O produto vetorial \(\vec{v}\times\vec{w}\)
  \pause
  \begin{itemize}[<+->]
      \item é ortogonal a $\vec{v}$
      \item é ortogonal a $\vec{w}$
      \item seu sentido é determinado pela regra da mão direita
  \end{itemize}
\end{frame}

%------------------------------------------------------------------------------%
\begin{frame}{Regra da Mão Direita}
  \onslide<+->{O sentido do produto vetorial é determinado pela regra da mão direita}
  \only<+->{\addgraphic{6}{9,2.5}{Right_hand_rule_cross_product}}

  \bigskip
  \begin{itemize}[<+->]
      \item aponte o dedo indicador na direção de $\vec{v}$
      \item aponte o dedo médio na direção a $\vec{w}$
      \item o polegar indica a direção de $\vec{v}\times\vec{w}$
  \end{itemize}

  \bigskip
  \onslide<+->{\emph{A ordem dos vetores é importante}}
  \onslide<+->{
  \[
    \vec{v}\times\vec{w} \;=\; -\left(\vec{w}\times\vec{v}\right)
  \]}
\end{frame}

%------------------------------------------------------------------------------%
\section{Definição}

%------------------------------------------------------------------------------%
\begin{frame}{Definição Algébrica}
  Dados \quad
  \(
    \vec{v} = \vetortri{v_1}{v_2}{v_3}
    \qquad
    \vec{w} = \vetortri{w_1}{w_2}{w_3}
  \)

  \pause
  \bigskip
  O produto vetorial é definido por
  \[
    \vec{v}\times\vec{w}
    = \vetortri{v_2w_3-v_3w_2}{v_3w_1-v_1w_3}{v_1w_2-v_2w_1}
  \]
\end{frame}

%------------------------------------------------------------------------------%
\begin{frame}{Representação Alternativa}
  Definindo os vetores
  \(
    \quad
    \vec{i} = \vetortri{1}{0}{0}
    \qquad
    \vec{j} = \vetortri{0}{1}{0}
    \qquad
    \vec{k} = \vetortri{0}{0}{1}
  \)

  \pause
  \bigskip
  Temos
   \[
    \vec{v}\times\vec{w}
    \;=\; \left( v_2 w_3 - v_3 w_2 \right) \vec{i}
    \;+\; \left( v_3 w_1 - v_1 w_3 \right) \vec{j}
    \;+\; \left( v_1 w_2 - v_2 w_1 \right) \vec{k}
  \]
\end{frame}

%------------------------------------------------------------------------------%
\begin{frame}{Forma por Determinante}
  Cálculo pelo determinante
  \[
    \vec{v}\times\vec{w}
    \;=\;
    \begin{vmatrix}[c]
      \vec{i} & \vec{j} & \vec{k} \\
      v_1     & v_2     & v_3     \\
      w_1     & w_2     & w_3
    \end{vmatrix}
  \]

  \pause
  Expandindo pela primeira linha
  \[
    \vec{v} \times \vec{w}
    \;=\;
    \begin{vmatrix}
      v_2 & v_3 \\
      w_2 & w_3
    \end{vmatrix}
    \vec{i}
    \;-\;
    \begin{vmatrix}
      v_1 & v_3 \\
      w_1 & w_3
    \end{vmatrix}
    \vec{j}
    \;+\;
    \begin{vmatrix}
      v_1 & v_2 \\
      w_1 & w_2
    \end{vmatrix}
    \vec{k}
  \]
  \pause
  Lembre-se do Teorema de Laplace -- \emph{Atenção aos sinais}
\end{frame}

%------------------------------------------------------------------------------%
\section{Exemplos}

%------------------------------------------------------------------------------%
\begin{question}
  Calcule o produto escalar e o produto vetorial do vetor $\vec{v}$ por $\vec{w}$
  \[
    \vec{v} = \vetortri{1}{2}{3}
    \qquad
    \vec{w} = \vetortri{2}{-1}{1}
  \]
\end{question}

%------------------------------------------------------------------------------%
\begin{resolution}{Produto Escalar}
  \begin{align*}
    \onslide<+->{\vec{v} \cdot \vec{w}}
    \onslide<+->{& = \vetortri{1}{2}{3} \cdot \vetortri{2}{-1}{1} \\}
    \onslide<+->{& = 1 \times 2 + 2(-1) + 3 \times 1 \\}
    \onslide<+->{& = 2 - 2 + 3 \\}
    \onslide<+->{& = 3}
  \end{align*}
\end{resolution}

%------------------------------------------------------------------------------%
\begin{resolution}{Produto Vetorial}
  \begin{align*}
    \onslide<+->{\vec{v} \times \vec{w}}
    \onslide<+->{& =   \vetortri{1}{2}{3} \times \vetortri{2}{-1}{1} }
    \onslide<+->{
    \;=\; \begin{vmatrix}
          \vec{i} & \vec{j} & \vec{k} \\
          1       & 2       & 3       \\
          2       & -1      & 1
        \end{vmatrix}}
    \\
    \onslide<+->{
    & = \begin{vmatrix}
           2 & 3 \\
          -1 & 1
        \end{vmatrix}
        \vec{i}}
    \onslide<+->{\;-\;
        \begin{vmatrix}
          1 & 3 \\
          2 & 1
        \end{vmatrix}
        \vec{j}}
    \onslide<+->{\;+\;
        \begin{vmatrix}
          1 & 2 \\
          2 & -1
        \end{vmatrix}
        \vec{k}}
    \\
    \onslide<+->{& = \big( 2 \times 1 - 3(-1)      \big) \vec{i}}
    \onslide<+->{  - \big( 1 \times 1 - 3 \times 2 \big) \vec{j}}
    \onslide<+->{  + \big( 1(-1)      - 2 \times 2 \big) \vec{k}}
    \\
    \onslide<+->{& =   5\vec{i} + 5\vec{j} - 5\vec{k}}
    \onslide<+->{\;=\; 5 \vetortri{1}{0}{0} + 5 \vetortri{0}{1}{0} -5 \vetortri{0}{0}{1}}
    \onslide<+->{\;=\; \vetortri{5}{5}{-5}}
  \end{align*}
\end{resolution}

%------------------------------------------------------------------------------%
\endinput

%------------------------------------------------------------------------------%
\begin{question}
  Dados os vetores
  \[
    \vec{v} = \vetortri{2}{-1}{4}
    \qquad
    \vec{w} = \vetortri{1}{3}{-2}
  \]
  Calcule \(\vec{v} \times \vec{w}\) e \(\vec{w} \times \vec{v}\)
\end{question}

%------------------------------------------------------------------------------%
\begin{resolution}{Produto Vetorial \(\vec{v} \times \vec{w}\)}
\begin{align*}
  \onslide<+->{\vec{v}\times\vec{w}}
  \onslide<+->{& = \begin{vmatrix}
    \vec{i} & \vec{j} & \vec{k} \\
    2       & -1      & 4       \\
    1       & 3       & -2
  \end{vmatrix}}
  \\
  \onslide<+->{& = \begin{vmatrix}
        -1 & 4 \\
         3 & -2
      \end{vmatrix}
      \vec{i}
      \;-\;
      \begin{vmatrix}
        2 & 4 \\
        1 & -2
      \end{vmatrix}
      \vec{j}
      \;+\;
      \begin{vmatrix}
        2 & -1 \\
        1 & 3
      \end{vmatrix}
      \vec{k}}
  \\
  \onslide<+->{
  & = \big( -1(-2) - 4 \times 3 \big) \vec{i}
    - \big( 2(-2) - 4  \times 1 \big) \vec{j}
    + \big( 2  \times 3 - (-1)1 \big) \vec{k}}
  \\
  \onslide<+->{& = -10 \vec{i} + 8 \vec{j} + 7 \vec{k}}
  \onslide<+->{\;=\; \vetortri{-10}{8}{7}}
\end{align*}
\end{resolution}

%------------------------------------------------------------------------------%
\begin{resolution}{Produto Vetorial \(\vec{w} \times \vec{v}\)}
\begin{align*}
  \onslide<+->{\vec{w}\times\vec{v}}
  \onslide<+->{& = \begin{vmatrix}
    \vec{i} & \vec{j} & \vec{k} \\
    1       & 3       & -2      \\
    2       & -1      &  4
  \end{vmatrix}}
  \\
  \onslide<+->{& = \begin{vmatrix}
       3 & -2 \\
      -1 & 4
    \end{vmatrix}
    \vec{i}
    \;-\;
    \begin{vmatrix}
      1 & -2 \\
      2 & 4
    \end{vmatrix}
    \vec{j}
    \;+\;
    \begin{vmatrix}
      1 & 3 \\
      2 & -1
    \end{vmatrix}
    \vec{k}}
  \\
  \onslide<+->{
  & = \big( 3 \times 4 - (-2)(-1) \big) \vec{i}
    - \big( 1 \times 4 - (-2)2 \big) \vec{j}
    + \big( 1(-1) - 3 \times 2 \big) \vec{k}}
  \\
  \onslide<+->{& = 10 \vec{i} - 8 \vec{j} - 7 \vec{k}}
  \onslide<+->{\;=\; \vetortri{10}{-8}{-7}}
\end{align*}
\end{resolution}

%------------------------------------------------------------------------------%
\endinput

%------------------------------------------------------------------------------%
\begin{question}
  Calcule o produto vetorial
  \[
    \vec{v} = \vetortri{1}{2}{-1}
    \qquad
    \vec{w} = \vetortri{2}{-1}{3}
  \]
  Em seguida verifique que o resultado é ortogonal aos dois vetores
\end{question}

%------------------------------------------------------------------------------%
\begin{resolution}{Produto Vetorial}
\begin{align*}
  \onslide<+->{\vec{\psi}}
  \onslide<+->{& = \vec{v}\times\vec{w}}
  \onslide<+->{\;=\; \begin{vmatrix}
    \vec{i} & \vec{j} & \vec{k} \\
    1       &  2      & -1      \\
    2       & -1      &  3
  \end{vmatrix}}
  \\
  \onslide<+->{& = \begin{vmatrix}
       2 & -1 \\
      -1 &  3
    \end{vmatrix}
    \vec{i}
    \;-\;
    \begin{vmatrix}
      1 & -1 \\
      2 &  3
    \end{vmatrix}
    \vec{j}
    \;+\;
    \begin{vmatrix}
      1 &  2 \\
      2 & -1
    \end{vmatrix}
    \vec{k}}
  \\
  \onslide<+->{
  & = \big( 2 \times 3 - (-1)(-1) \big) \vec{i}
    - \big( 1 \times 3 - (-1)2 \big) \vec{j}
    + \big( 1(-1) - 2 \times 2 \big) \vec{k}}
  \\
  \onslide<+->{& = 5 \vec{i} -5 \vec{j} -5 \vec{k}}
  \onslide<+->{\;=\; \vetortri{5}{-5}{-5}}
\end{align*}
\end{resolution}

%------------------------------------------------------------------------------%
\begin{resolution}{Verificando a Ortogonalidade}
\begin{columns}
\column{0.5\textwidth}
  \begin{align*}
    \onslide<+->{\vec{\psi} \cdot \vec{v}}
    \onslide<+->{& = \vetortri{5}{-5}{-5} \cdot \vetortri{1}{2}{-1} \\}
    \onslide<+->{& = 5 \times 1 + (-5)2 + (-5)(-1) \\}
    \onslide<+->{& = 5 - 10 + 5 \\}
    \onslide<+->{& = 0}
  \end{align*}
  \onslide<+->{\(\vec{\psi}\) e \(\vec{v}\) são ortogonais}
\column{0.5\textwidth}
  \begin{align*}
    \onslide<+->{\vec{\psi} \cdot \vec{w}}
    \onslide<+->{& = \vetortri{5}{-5}{-5} \cdot \vetortri{2}{-1}{3} \\}
    \onslide<+->{& = 5 \times 2 + (-5)(-1) + (-5)3 \\}
    \onslide<+->{& = 10 + 5 - 15 \\}
    \onslide<+->{& = 0}
  \end{align*}
  \onslide<+->{\(\vec{\psi}\) e \(\vec{w}\) são ortogonais}
\end{columns}
\end{resolution}

%------------------------------------------------------------------------------%
\endinput

%------------------------------------------------------------------------------%
\begin{question}
  Determine um vetor unitário ortogonal aos vetores
  \[
    \vec{v} = \vetortri{2}{-1}{1}
    \qquad
    \vec{w} = \vetortri{-1}{3}{2}
  \]

  \pause
  \bigskip
  Usamos o produto vetorial para encontrar um vetor ortogonal e \\
  dividimos ele por sua norma para encontrar um vetor unitário
\end{question}

%------------------------------------------------------------------------------%
\begin{resolution}{Produto Vetorial}
\begin{align*}
  \onslide<+->{\vec{\psi}}
  \onslide<+->{& = \vec{v}\times\vec{w}}
  \onslide<+->{\;=\; \begin{vmatrix}
    \vec{i} & \vec{j} & \vec{k} \\
     2      & -1      &  1      \\
    -1      &  3      &  2
  \end{vmatrix}}
  \\
  \onslide<+->{& = \begin{vmatrix}
      -1 &  1 \\
       3 &  2
    \end{vmatrix}
    \vec{i}
    \;-\;
    \begin{vmatrix}
       2 &  1 \\
      -1 &  2
    \end{vmatrix}
    \vec{j}
    \;+\;
    \begin{vmatrix}
       2 &  -1 \\
      -1 &   3
    \end{vmatrix}
    \vec{k}}
  \\
  \onslide<+->{
  & = \big( (-1)2 - 1 \times 3 \big) \vec{i}
    - \big( 2 \times 2 - 1(-1) \big) \vec{j}
    + \big( 2 \times 3 - (-1)(-1) \big) \vec{k}}
  \\
  \onslide<+->{& = -5 \vec{i} -5 \vec{j} + 5 \vec{k}}
  \onslide<+->{\;=\; \vetortri{-5}{-5}{5}}
\end{align*}
\end{resolution}

%------------------------------------------------------------------------------%
\begin{resolution}{Vetor Unitário}
\begin{columns}[t]
\column{0.3\textwidth}
  \onslide<+->{Norma de \(\vec{\psi}\)}
  \begin{align*}
    \onslide<+->{\norm{\vec{\psi}}}
    \onslide<+->{& = \norm{\vetortri{-5}{-5}{5}} \\}
    \onslide<+->{& = \sqrt{(-5)^2 + (-5)^2 + 5^2} \\}
    \onslide<+->{& = \sqrt{3 \times 5^2} \\}
    \onslide<+->{& = 5\sqrt{3}}
  \end{align*}
\column{0.5\textwidth}
  \onslide<+->{Vetor unitário da direção e sentido de \(\vec{\psi}\)}
  \begin{align*}
    \onslide<+->{\hat{\psi}}
    \onslide<+->{& = \frac{\vec{\psi}}{\norm{\vec{\psi}}} }
    \onslide<+->{\;=\; \frac{1}{5\sqrt{3}} \vetortri{-5}{-5}{5} \\}
    \onslide<+->{& = \Vetortri{\frac{-1}{\sqrt{3}}}{\frac{-1}{\sqrt{3}}}{\frac{1}{\sqrt{3}}}}
  \end{align*}
\end{columns}
\end{resolution}

%------------------------------------------------------------------------------%
\endinput

%------------------------------------------------------------------------------%
\begin{question}
  Determine todos os vetores simultaneamente perpendiculares
  a $\vec{v}$ e $\vec{w}$ \\
  e que possuem módulo 10
  \[
    \vec{v} = \vetortri{1}{2}{3}
    \qquad
    \vec{w} = \vetortri{2}{-1}{1}
  \]
\end{question}

%------------------------------------------------------------------------------%
\begin{resolution}{Produto Vetorial}
\begin{align*}
  \onslide<+->{\vec{\psi}}
  \onslide<+->{& = \vec{v}\times\vec{w}}
  \onslide<+->{\;=\; \begin{vmatrix}
    \vec{i} & \vec{j} & \vec{k} \\
     1      &   2     &  3      \\
     2      &  -1     &  1
  \end{vmatrix}}
  \\
  \onslide<+->{& = \begin{vmatrix}
       2 &  3 \\
      -1 &  1
    \end{vmatrix}
    \vec{i}
    \;-\;
    \begin{vmatrix}
       1 &  3 \\
       2 &  1
    \end{vmatrix}
    \vec{j}
    \;+\;
    \begin{vmatrix}
       1 &  2 \\
       2 & -1
    \end{vmatrix}
    \vec{k}}
  \\
  \onslide<+->{
  & = \big( 2 \times 1 - 3(-1) \big) \vec{i}
    - \big( 1 \times 1 - 3 \times 2 \big) \vec{j}
    + \big( 1(-1) - 2 \times 2 \big) \vec{k}}
  \\
  \onslide<+->{& = 5 \vec{i} +5 \vec{j} - 5 \vec{k}}
  \onslide<+->{\;=\; \vetortri{5}{5}{-5}}
\end{align*}
\end{resolution}

%------------------------------------------------------------------------------%
\begin{resolution}{Vetor Unitário}
\begin{columns}[t]
\column{0.3\textwidth}
  \onslide<+->{Norma de \(\vec{\psi}\)}
  \begin{align*}
    \onslide<+->{\norm{\vec{\psi}}}
    \onslide<+->{& = \norm{\vetortri{5}{5}{-5}} \\}
    \onslide<+->{& = \sqrt{5^2 + 5^2 + (-5)^2} \\}
    \onslide<+->{& = \sqrt{3\times 5^2} \\}
    \onslide<+->{& = 5\sqrt{3}}
  \end{align*}
\column{0.5\textwidth}
  \onslide<+->{Vetor unitário da direção e sentido de \(\vec{\psi}\)}
  \begin{align*}
    \onslide<+->{\hat{\psi}}
    \onslide<+->{& = \frac{\vec{\psi}}{\norm{\vec{\psi}}} }
    \onslide<+->{\;=\; \frac{1}{5\sqrt{3}} \vetortri{5}{5}{-5} \\}
    \onslide<+->{& = \Vetortri{\frac{1}{\sqrt{3}}}{\frac{1}{\sqrt{3}}}{\frac{-1}{\sqrt{3}}}}
  \end{align*}
\end{columns}
\end{resolution}

%------------------------------------------------------------------------------%
\begin{resolution}{Vetores com Módulo 10}
  Temos dois vetores na direção de \(\vec{\psi}\)

  \pause
  que são ortogonais a \(\vec{v}\) e \(\vec{w}\)
  \[
    \pause
    \vec{a} = 10 \hat{\psi}
  \]
  \[
    \pause
    \vec{b} = -10 \hat{\psi}
  \]
\end{resolution}

%------------------------------------------------------------------------------%
\begin{resolution}{Vetores com Módulo 10}
  \[
    \vec{a}
    \pause
    \;=\; 10 \hat{\psi}
    \pause
    \;=\; 10 \Vetortri{\frac{5}{\sqrt{59}}}{\frac{-5}{\sqrt{59}}}{\frac{-3}{\sqrt{59}}}
    \pause
    \;=\; \Vetortri{\frac{50}{\sqrt{59}}}{\frac{-50}{\sqrt{59}}}{\frac{-30}{\sqrt{59}}}
  \]
\end{resolution}

%------------------------------------------------------------------------------%
\begin{resolution}{Vetores com Módulo 10}
  \[
    \vec{b}
    \pause
    \;=\; -10 \hat{\psi}
    \pause
    \;=\; -10 \Vetortri{\frac{5}{\sqrt{59}}}{\frac{-5}{\sqrt{59}}}{\frac{-3}{\sqrt{59}}}
    \pause
    \;=\; \Vetortri{\frac{-50}{\sqrt{59}}}{\frac{50}{\sqrt{59}}}{\frac{30}{\sqrt{59}}}
  \]
\end{resolution}

%------------------------------------------------------------------------------%
\endinput


%%------------------------------------------------------------------------------%
\begin{question}
  Uma força é aplicada em uma estrutura segundo

  \[
  \vec F=(4,-2,5)
  \]

  e o braço de alavanca é representado por

  \[
  \vec r=(2,1,-1)
  \]

  O momento da força em relação à origem é

  \[
  \vec M=\vec r\times\vec F
  \]

  Determine $\vec M$ e seu módulo
\end{question}

%------------------------------------------------------------------------------%
\begin{resolution}{Produto Vetorial}

\end{resolution}

%------------------------------------------------------------------------------%
\endinput

%%------------------------------------------------------------------------------%
\begin{question}
  Considere os pontos

  \[
  A=(1,2,0)
  \]

  \[
  B=(4,-1,2)
  \]

  \[
  C=(2,3,5)
  \]

  Determine um vetor perpendicular ao plano determinado por $A$, $B$ e $C$
\end{question}

%------------------------------------------------------------------------------%
\begin{resolution}{Produto Vetorial}

\end{resolution}

%------------------------------------------------------------------------------%
\endinput


%%------------------------------------------------------------------------------%
%\section{Lista Mínima}
%
%%------------------------------------------------------------------------------%
%\begin{frame}{Lista Mínima}
%  \begin{enumerate}
%    \item
%  \end{enumerate}
%
%  \vfill
%  Atenção: A prova é baseada no livro, não nas apresentações
%\end{frame}

%------------------------------------------------------------------------------%
\end{document}
%------------------------------------------------------------------------------%
